% =====================================================================
%  CHƯƠNG 7. THỰC NGHIỆM LẬP TRÌNH TÁC TỬ HỌC TĂNG CƯỜNG TRÊN JUPYTER NOTEBOOK
% =====================================================================
\chapter{Thực nghiệm lập trình tác tử học tăng cường trên Jupyter Notebook}\label{ch:thucnghiem}

\section{Mô tả bài toán điều hướng mê cung}
Để hiện thực hóa cơ chế tự học của Learning Agent (Chương~4) bằng thuật toán Q-Learning (Chương~5), bài thực hành xây
dựng chương trình hoàn chỉnh trên Jupyter Notebook \texttt{agent01\_rl\_qlearning.ipynb} giải bài toán \textbf{robot
điều hướng mê cung} theo tài liệu GeeksforGeeks \cite{gfg_rl}. Mê cung do sinh viên tự thiết kế, kích thước
$\SL{ch8_qlearning/params/n_rows} \times \SL{ch8_qlearning/params/n_cols}$ với \SL{ch8_qlearning/params/n_fire} ô tường lửa
(Fire). Xuất phát từ ô $S\,\SL{ch8_qlearning/params/start}$, robot không biết trước bản đồ; nó phải tự khám phá, học cách
tránh tường lửa và tìm ra lộ trình ngắn nhất về đích $G\,\SL{ch8_qlearning/params/goal}$ với tổng phần thưởng lớn nhất
(Hình~\ref{fig:ch7-maze}).

\hinhtikz{ch7_maze}{Mê cung $12 \times 12$ của bài thực hành: ô xuất phát $S$, ô đích $G$, các ô tường lửa, bốn hành
động và hàm phần thưởng.}{ch7-maze}

\subsection{Bốn thành tố của bài toán học tăng cường}
Bài toán được mô hình hóa thành một MDP (Bảng~\ref{tab:ch7-mdp}):
\begin{enumerate}
  \item \textbf{Trạng thái} $\mathcal{S}$: tọa độ ô $(r, c)$ với $0 \le r, c \le 11$, tổng cộng
        \SL{ch8_qlearning/params/n_states} trạng thái lưới. Mỗi trạng thái được đánh chỉ số $r \cdot 12 + c$ để làm
        hàng của Q-Table.
  \item \textbf{Hành động} $\mathcal{A}$: bốn hướng rời rạc Trái $(0,-1)$, Phải $(0,1)$, Lên $(-1,0)$, Xuống $(1,0)$.
  \item \textbf{Phần thưởng} $R$: $R_{\text{fire}} = \SL{ch8_qlearning/params/r_fire}$ khi đâm vào tường lửa hoặc ra
        khỏi mê cung (robot đứng yên tại chỗ), $R_{\text{goal}} = +\SL{ch8_qlearning/params/r_goal}$ khi về đích (kết
        thúc episode), $R_{\text{step}} = \SL{ch8_qlearning/params/r_step}$ cho mỗi bước đi hợp lệ để thúc tác tử đi
        đường ngắn.
  \item \textbf{Quy tắc học}: cập nhật Bellman $Q(s,a) \leftarrow Q(s,a) + \alpha\,[R + \gamma \max_{a'} Q(s',a') -
        Q(s,a)]$ với $\alpha = \SL{ch8_qlearning/params/alpha}$, $\gamma = \SL{ch8_qlearning/params/gamma}$ và lịch
        $\varepsilon$-greedy giảm từ \SL{ch8_qlearning/params/eps_start} về \SL{ch8_qlearning/params/eps_min} (nhân
        \SL{ch8_qlearning/params/eps_decay} sau mỗi episode), huấn luyện \SL{ch8_qlearning/params/episodes} episode,
        tối đa \SL{ch8_qlearning/params/max_steps} bước mỗi episode.
\end{enumerate}

\begin{table}[!htbp]\centering
\caption{Các thành phần MDP của bài toán mê cung.}\label{tab:ch7-mdp}
\small
\begin{tabularx}{\linewidth}{L{3.4cm}X}
\toprule
\textbf{Thành phần} & \textbf{Giá trị trong bài thực hành} \\
\midrule
$\mathcal{S}$ & \SL{ch8_qlearning/params/n_states} ô $(r, c)$; robot chỉ đứng được trên ô trống \\
$\mathcal{A}$ & \{Trái, Phải, Lên, Xuống\} \\
$P(s' \mid s, a)$ & Tất định: sang ô kề theo hướng $a$; nếu ô đó là tường lửa hoặc ngoài biên thì $s' = s$ \\
$R(s, a, s')$ & $+\SL{ch8_qlearning/params/r_goal}$ về đích, $\SL{ch8_qlearning/params/r_fire}$ đâm tường/biên,
$\SL{ch8_qlearning/params/r_step}$ bước thường \\
$\gamma$ & \SL{ch8_qlearning/params/gamma} \\
Trạng thái kết thúc & Ô đích $G\,\SL{ch8_qlearning/params/goal}$, hoặc hết \SL{ch8_qlearning/params/max_steps} bước \\
\bottomrule
\end{tabularx}
\end{table}

\section{Môi trường lập trình}
Notebook chạy trên JupyterLab~4 với Python~3.13, dùng NumPy cho Q-Table, pandas để hiển thị bảng và Matplotlib để vẽ.
Q-Table chỉ có $\SL{ch8_qlearning/params/n_states} \times 4$ ô số thực nên toàn bộ việc huấn luyện chạy trên CPU và mất
\SL{ch8_qlearning/result/seconds} giây cho \SL{ch8_qlearning/params/episodes} episode; bài toán dạng bảng này không
cần GPU. Notebook được tổ chức thành tám phần (Hình~\ref{fig:nb-window}):
\begin{enumerate}
  \item Thư viện; \quad 2.~Môi trường mê cung; \quad 3.~Q-Learning Agent và vòng huấn luyện; \quad 4.~Huấn luyện;
  \setcounter{enumi}{4}
  \item Lộ trình tối ưu theo chính sách tham lam; \quad 6.~Trực quan hóa; \quad 7.~Trích Q-Table;
        \quad 8.~Lưu số liệu.
\end{enumerate}

\anhNB{nb_window}{Ảnh chụp màn hình JupyterLab mở notebook Q-Learning của bài tập: phần đầu notebook và ô mã thư
viện.}{nb-window}

\section{Ảnh chụp màn hình mã nguồn thuật toán (Copy Screen Code)}
Các hình dưới đây là ảnh chụp màn hình thực tế những khối mã xây dựng môi trường, thiết lập tham số và thuật toán
Q-Learning Agent trên Jupyter Notebook, đúng yêu cầu của bài tập (copy screen code). Sau mỗi ảnh chụp, báo cáo in lại
đoạn mã tương ứng ở dạng văn bản (trích nguyên văn từ mã đã chạy) kèm giải thích.

\subsection{Khối thư viện}
\anhNB[0.55]{code_imports}{Ảnh chụp màn hình ô mã 1: nạp thư viện.}{nb-imports}

\subsection{Khối môi trường mê cung}
\anhNB{code_maze_1}{Ảnh chụp màn hình ô mã 2 (phần 1/2): bản đồ mê cung, hành động, phần thưởng và đầu lớp
\texttt{MazeEnv}.}{nb-maze}
\anhNB{code_maze_2}{Ảnh chụp màn hình ô mã 2 (phần 2/2): phương thức \texttt{step} và BFS kiểm chứng của lớp
\texttt{MazeEnv}.}{nb-maze-2}

Bản đồ được khai báo bằng các chuỗi ký tự để dễ đọc và dễ sửa: \texttt{S} là xuất phát, \texttt{G} là đích,
\texttt{\#} là tường lửa.

\maNguon{maze__LAYOUT}{Bản đồ mê cung $12 \times 12$.}

Lớp \texttt{MazeEnv} cài đặt hàm chuyển trạng thái $P$ và hàm phần thưởng $R$ của MDP. Phương thức \texttt{step} nhận
một hành động, tính ô đích đến; nếu ô đó ra khỏi mê cung hoặc là tường lửa thì robot đứng yên và nhận
$R_{\text{fire}}$; nếu là ô đích thì episode kết thúc với $R_{\text{goal}}$; còn lại nhận $R_{\text{step}}$. Phương
thức \texttt{shortest\_path\_length} chạy BFS trên bản đồ. Tác tử \emph{không} dùng hàm này; nó chỉ để kiểm chứng sau
huấn luyện rằng lộ trình học được có đúng là ngắn nhất.

\maDai{maze__MazeEnv}{Lớp môi trường \texttt{MazeEnv}: trạng thái, bước chuyển, phần thưởng và BFS kiểm chứng.}

\subsection{Khối Q-Learning Agent và vòng huấn luyện}
\anhNB{code_qlearning_1}{Ảnh chụp màn hình ô mã 3 (phần 1/2): siêu tham số và lớp \texttt{QLearningAgent}.}{nb-qlearning}
\anhNB{code_qlearning_2}{Ảnh chụp màn hình ô mã 3 (phần 2/2): hàm \texttt{train} và
\texttt{greedy\_path}.}{nb-qlearning-2}

Lớp \texttt{QLearningAgent} giữ Q-Table dạng mảng NumPy kích thước $\SL{ch8_qlearning/params/n_states} \times 4$,
khởi tạo bằng 0. Phương thức \texttt{choose\_action} hiện thực chính sách $\varepsilon$-greedy; khi nhiều hành động có
cùng giá trị lớn nhất (đầu quá trình mọi giá trị đều bằng 0), tác tử chọn ngẫu nhiên giữa chúng để không thiên vị
hướng nào. Phương thức \texttt{update} là đúng một dòng của phương trình Bellman; ở trạng thái kết thúc, mục tiêu TD chỉ
là phần thưởng $r$.

\maDai{qlearning__QLearningAgent}{Lớp \texttt{QLearningAgent}: Q-Table, chính sách $\varepsilon$-greedy, cập nhật
Bellman và giảm $\varepsilon$.}

Hàm \texttt{train} lặp qua các episode. Mỗi episode bắt đầu tại $S$; ở mỗi bước tác tử chọn hành động, môi trường trả
về trạng thái mới và phần thưởng, tác tử cập nhật Q-Table. Episode dừng khi về đích hoặc khi chạm giới hạn
\texttt{max\_steps}. Sau mỗi episode, $\varepsilon$ giảm một bậc; cứ 250 episode hàm in một dòng nhật ký gồm
$\varepsilon$, phần thưởng trung bình, số bước trung bình và tỉ lệ về đích.

\maDai{qlearning__train}{Hàm \texttt{train}: vòng huấn luyện qua các episode và nhật ký tiến trình.}

Sau huấn luyện, \texttt{greedy\_path} chạy chính sách tham lam $\pi^*(s) = \arg\max_a Q(s,a)$ từ $S$, không còn thăm dò,
để lấy lộ trình tối ưu mà tác tử đã học.

\maNguon{qlearning__greedy_path}{Hàm \texttt{greedy\_path}: lộ trình theo chính sách tham lam.}

\subsection{Khối trực quan hóa và trích Q-Table}
\anhNB{code_viz}{Ảnh chụp màn hình ô mã 6: các hàm vẽ mê cung, đường cong phần thưởng và trích Q-Table.}{nb-viz}

Hàm \texttt{draw\_maze} tô tường lửa, vẽ lộ trình lên lưới; \texttt{draw\_rewards} vẽ phần thưởng từng episode, đường
trung bình trượt 100 episode và lịch $\varepsilon$ trên trục phụ. Hàm \texttt{q\_table\_excerpt} tự chọn các ô then
chốt của lộ trình: ô xuất phát, mọi \emph{điểm rẽ} (nơi hướng đi thay đổi) và ô ngay trước đích.

\maNguon{viz__q_table_excerpt}{Hàm \texttt{q\_table\_excerpt}: trích Q-Table tại các ô then chốt của lộ trình.}
\maDai{viz__draw_maze}{Hàm \texttt{draw\_maze}: vẽ mê cung và lộ trình tối ưu.}
\maDai{viz__draw_rewards}{Hàm \texttt{draw\_rewards}: đường cong hội tụ phần thưởng và lịch $\varepsilon$.}
